\subsection{The family maximiser}
\label{sec:entropy-local-seven}

Write \(E_3(z)=35^{-1}\sum_{|S|=3}S(\rho_S)\) for a unit
seven-qubit vector, and retain the two real amplitude patterns
\[
 U=\sum_{i<j}s_{ij}\bigl(|0,e_i+e_j\rangle
       -|1,\mathbf1-e_i-e_j\rangle\bigr),\qquad
 V=|0,\mathbf1\rangle+|1,\mathbf0\rangle,
\]
so \(\psi_x=\alpha U+\beta V\), \(\alpha=\sqrt{(1-2x)/30}\), \(\beta=\sqrt x\),
\(\langle U,V\rangle=0\), \(\|U\|^2=30\), and \(\|V\|^2=2\).

\begin{lemma}[Symmetric criticality of the conference family; proved, ES-1]
\label{lem:seven-family-symmetric-criticality}
Let $\mathcal F$ be a real function on the unit sphere, invariant under
local unitaries, qubit permutations and complex conjugation, and
differentiable at $\psi_x$ for some $0<x<1/2$. If $x$ is a stationary point
of $x\mapsto\mathcal F(\psi_x)$, then $\psi_x$ is stationary on the whole
unit sphere. This applies to $\pi_3$ (a polynomial) and to $E_3$, whose
marginals have full rank for $0<x<1/2$ by Theorem~\ref{thm:ico}.
\end{lemma}
\begin{proof}
We exhibit a common symmetry group of $U,V$ whose real fixed space is
exactly $\operatorname{span}_{\mathbb R}\{U,V\}$. Ordinary labels are $0,\ldots,5$,
with the distinguished qubit denoted by $s$.
The diagonal unitary
\[
 T=-Z_s\otimes\operatorname{diag}(1,i)^{\otimes6}
\]
fixes $U,V$. Its eigenvalue on $|s,b\rangle$ is $-(-1)^s i^{|b|}$,
so its fixed space consists of the pair and complement sectors
$(s,|b|)=(0,2),(1,4)$ and the two endpoints $(0,6),(1,0)$.

Put $g=(1\,3\,2\,4\,5)$, $p=(0\,1)(2\,4)$ and
$\epsilon=(1,1,1,-1,1,-1)$. The conference signs satisfy
$s_{g(i)g(j)}=s_{ij}$ and
$\epsilon_i\epsilon_j s_{p(i)p(j)}=s_{ij}$.
If $P_\pi$ permutes the ordinary qubits by $\pi$, then $P_g$ and
\[
 W=\left(-Z_s\otimes\bigotimes_{i=0}^5
           \operatorname{diag}(1,-i\epsilon_i)\right)P_p
\]
fix $U,V$ and both endpoints. On either pair sector, dividing the
coefficients by $s_{ij}$ turns their actions into ordinary permutations
of the fifteen pairs. These permutations act transitively: $g$ has
three pair orbits (pairs through $0$, pentagon edges, and pentagon
diagonals), while $p$ sends $\{0,3\}$ to $\{1,3\}$ and
$\{0,2\}$ to $\{1,4\}$, connecting all three.
Thus a real vector fixed by $T,P_g,W$ has one coefficient on each
signed pair sector and one on each endpoint.

Finally the local unitary
\[
 Q=\left(\operatorname{diag}(i,1)_s\otimes
       \operatorname{diag}(1,e^{i\pi/4})^{\otimes6}\right)X^{\otimes7}
\]
fixes $U,V$. On the displayed support it exchanges complementary
pair strings with a minus sign and exchanges the endpoints with a
plus sign. Its fixed space within the preceding four-dimensional
space is therefore exactly $\operatorname{span}_{\mathbb R}\{U,V\}$.

At a real $\psi_x$, conjugation invariance makes the spherical gradient
real. The four displayed symmetries act orthogonally on
$\mathbb C^{128}\cong\mathbb R^{256}$, preserve $\mathcal F$ and fix
$\psi_x$, so they fix that gradient, which therefore lies in this plane. Orthogonality to $\psi_x$ leaves only the
one-dimensional family tangent. Stationarity within the family makes that last component vanish,
proving full-space stationarity.
\end{proof}

\begin{proposition}[Full-space stationarity; proved, ES-1]
\label{prop:entropy-seven-stationarity}
The unique maximiser \(\psi_{x_*}\) of \(E_3\) within the icosahedral
family is a stationary point on the whole projective space
\(\mathbb CP^{127}\). Its parameter is the unique root in \((0,1/2)\) of
\begin{equation}\label{eq:entropy-seven-star}
 2^{32}(1-2x_*)^{15}=3^8(1+8x_*)^8(1+28x_*)^7,
 \qquad x_*=0.069828\ldots.
\end{equation}
In particular \(E_3(\psi_{x_*})=2.96535602947767\ldots\) bits, and
\(1/32<x_*<37/524\), so \(x_*/u_*\in(1,37/15)\) with \(u_*=(1-2x_*)/30\).
\end{proposition}
\begin{proof}
By Lemma~\ref{lem:seven-family-symmetric-criticality}, it suffices to
find the family stationary point. The marginal spectra of
Theorem~\ref{thm:ico} give
\[
 E_3'(x)=\frac4{35\ln2}\ln
  \frac{(4u)^{13}(8u)^2}{(3u+x)^8(u+x)^7},
 \qquad u=(1-2x)/30.
\]
Strict concavity and the endpoint derivative signs in
Corollary~\ref{cor:ico-entropy-separation} give a unique interior
zero, equivalent to~\eqref{eq:entropy-seven-star}. At this zero the symmetry
lemma gives full-projective stationarity. At \(x=1/32\) one has \(u=x\) and
\(E_3'(1/32)=4\ln512/(35\ln2)>0\), so \(1/32<x_*<37/524\) by concavity
and Corollary~\ref{cor:ico-entropy-separation}; since \(x/u\) increases
with \(x\) and equals \(37/15\) at \(x=37/524\), \(x_*/u_*\in(1,37/15)\).
\end{proof}

Throughout the entropy sections, \(\phi=\psi_{x_*}\), \(u_*=(1-2x_*)/30\) and
\(E_*=E_3(\phi)\); a \emph{competitor} is a pure state \(z\) with \(E_3(z)\ge E_*\).
